\documentclass{jsarticle}
\usepackage{amsmath}
\usepackage{amssymb}
\begin{document}
\title{}
\author{}
\date{}
\maketitle

メタン、エタン、プロパン、ブタン、ヘプタン、ヘキサン、
ヘプタン、オクタン、ノナン、デカン

付加基が多い方から、位置番号をつけて、

モノ、ジ、テトラ、ブタン、ペンタ、ヘキサ、ヘプタ、オクタ、ノナ、デカ

と、数字を表して、
メタノール、エタノール、プロパノール、ブタノール、
ナトリウムアルコキシド、ナトリウムフェノキシド、
１，２−メチルプロパノールが$CH_3-CH(CH_3)-CH_2-CH_2-OH$
と表せられて、
これらと、
ビリン化合物、コリン化合物、は、恐るべし構造式で表されていて、
フラバンが、安心する構造式であり、
ビタミン、ホルモン、糖類が、きれいな有機化合物の構造式で表せられていて、
これらによって、すべての命名法が示されている事典になっている。

\end{document}
